\medskip



Les parties I et II peuvent être abordées de façon indépendante.

\begin{center}\textbf{Partie I : Effet de l'introduction d'une nouvelle espèce.}\end{center}






Soit $f$ la fonction définie pour   $t \in [0~;~120]$ par:

\[f(t) = \left(0,04t^2 - 8t + 400\right)\text{e}^{\frac{t}{50}} + 40.\]

La variable $t$ représente le temps écoulé, en jour, à partir de l'introduction à l'instant $t = 0$ des truites dans le lac, et $f(t)$ modélise le nombre de crapauds à l'instant $t$.

\medskip

\begin{enumerate}
\item 
Pour $t = 0$,\, $f(0) = 400 \e^{0} + 40 = 440$~(crapauds)
\item 



Sur l'intervalle [0~;~120] en dérivant le produit :

$\begin{gathered}f'(t) = (0,08t - 8)\text{e}^{\frac{t}{50}} + \dfrac{1}{50}\left(0,04t^2 - 8t + 400\right)\text{e}^{\frac{t}{50}} = \text{e}^{\frac{t}{50}}\left[(0,08t - 8) + \dfrac{1}{50}\left(0,04t^2 - 8t + 400\right)\right]\\
\phantom{f'(t)} 
= \text{e}^{\frac{t}{50}}\left(0,08t - 8 + \np{0,0008}t^2 - 0,16 t + 8 \right)
= \text{e}^{\frac{t}{50}}\left(-0,08t + \np{0,0008}t^2\right) \\
\phantom{f'(t)} 
= \np{0,0008}\left(t^2 - 100t\right)\text{e}^{\frac{t}{50}} 
= t(t - 100)\text{e}^{\frac{t}{50}} \times 8 \times 10^{-4}.\end{gathered}$



\item 
On sait que quel que soit $t \in [0~;~120]$, $\text{e}^{\frac{t}{50}} > 0$ : le signe de $f'(t)$ est donc celui du trinôme $t(t - 100)$ qui est positif sauf sur l'intervalle $]0~;~100[$ (entre les racines du trinôme). 

$f(0) = 480$, $f(100) = (400 - 800 + 400)\text{e}^{\frac{100}{50}}  + 40 = 0 + 40 = 40$ et \\
$f(120) = 576  - 960 + 400)\text{e}^{\frac{120}{50}} + 40 = 16 \text{e}^{2,4} + 40 \approx 216,37$.

On dresse le tableau de variations de $f$ sur l'intervalle $[0~;~120]$:

\begin{center}
{\renewcommand{\arraystretch}{1.3}
\psset{nodesep=3pt,arrowsize=2pt 3}  
\def\esp{\hspace*{1.5cm}}
\def\hauteur{0pt}
$\begin{array}{|c|l *3{c} c|}
\hline
 t & 0 & \esp & 100 & \esp & 120 \\
 \hline
t(t-100) & 0 &  \pmb{-} & \vline\hspace{-2.7pt}0 & \pmb{+} & \\  
 \hline
f'(t) &  0 &  \pmb{-} & \vline\hspace{-2.7pt}0 & \pmb{+} & \\  
\hline
  & \Rnode{max1}{440}  &  &  &  & \Rnode{max2}{216,37}   \\
f & &  & & &  \rule{0pt}{\hauteur} \\
 &  & &   \Rnode{min}{40} & & \rule{0pt}{\hauteur}
\ncline{->}{max1}{min} \ncline{->}{min}{max2}\\
\hline
\end{array}$
}
\end{center}


 \item Selon cette modélisation:
	\begin{enumerate}
		\item 
On a vu dans la question précédente que $f(100) = 40$ est le minimum de la fonction $f$, on a donc $J = 100$.
		\item 
		On a aussi vu que de 100 jours à 120 jours le nombre de crapauds croît strictement de 40 à environ 216 : il dépassera donc 140 individus.
		\item 
		
Soit $J_{140}$ la durée en jour à partir de laquelle le nombre de crapauds dépassera $140$ individus. La calculatrice donne :
		
		$f(115) \approx 130<140$ et $f(116) \approx 144>140$, donc $J_{140} = 116$.
	\end{enumerate}
\end{enumerate}


\begin{center}\textbf{Partie II : Effet de la Chytridiomycose sur une population de têtards}\end{center}




















\medskip

\begin{enumerate}
\item On complète l'arbre de probabilité suivant en utilisant les données de l'énoncé:

\begin{center}
\pstree[treemode=R,nodesepB=3pt,levelsep=2.75cm]{\TR{}}
{\pstree{\TR{$L$~~}\taput{0,25}}
	{\TR{$T$}\taput{0,74}
	\TR{$\overline{T}$}\tbput{0,24}
	}
\pstree{\TR{$\overline{L}$~~}\tbput{0,75}}
	{\TR{$T$}\taput{0}
	\TR{$\overline{T}$}\tbput{1}
	}	
}
\end{center}

\item 
D'après la loi des probabilités totales : 

$P(T) = P(L \cap T) + P\left(\overline{L} \cap T\right) 
= P(L) \times P_L(T) + P\left(\overline{L}\right) \times P_{\overline{L}}(T) 
= 0,25 \times 0,74 + 0,75 \times 0 
= 0,185$

\item Le têtard n'est pas contaminé. La probabilité que le lac soit infecté est:

$P_{\overline{T}}(L) 
 = \dfrac{P\left(L \cap \overline{T}\right)}{P\left(\overline{T}\right)} = \dfrac{0,25 \times 0,26}{1 - 0,185} = \dfrac{0,065}{0,815}  \approx 0,0797$, soit $0,080$ au millième près.
\end{enumerate}

\medskip
